\documentclass[11pt,a4paper]{article}
\usepackage[margin=19mm,headheight=15pt]{geometry}
\usepackage{fontspec}
\setmainfont{Noto Serif}
\usepackage{amsmath,amssymb,mathtools}
\usepackage{xcolor}
\usepackage{fancyhdr}
\usepackage{enumitem}
\usepackage{microtype}
\setlength{\parindent}{0pt}
\setlength{\parskip}{6pt}
\definecolor{navy}{HTML}{183153}
\definecolor{lightgray}{HTML}{F3F5F7}
\pagestyle{fancy}\fancyhf{}
\lhead{\small ME 5309 -- Mid-Semester Examination}
\rhead{\small ME 5309}
\cfoot{\thepage}
\begin{document}
\begin{center}
{\Large\bfseries ME 5309 -- Mid-Semester Examination}\\[3pt]
{\large Answer Booklet}\\[8pt]
\end{center}
\textbf{Instructions:} Attempt any three questions. Each question carries 10 marks. Show working where appropriate.
\section*{Question 1 \hfill [10 marks]}
\subsection*{Q1(a) \hfill [4 marks]}

\vspace{2mm}

\subsection*{Q1(b) \hfill [3 marks]}

ADC step = 5/2\textasciicircum{}12 = 0.0012207 V = 1.2207 mV.
Sensor sensitivity is 10 mV/N, so ΔF = 1.2207/10 = 0.12207 N.
The smallest resolvable force change is approximately 0.122 N.\par\smallskip\textit{Calculation:} $2^{12}=4096$ levels. One step is $5/4096=0.0012207$ V, i.e. $1.2207$ mV. Dividing by the sensitivity, $1.2207\,\mathrm{mV}/(10\,\mathrm{mV/N})=0.12207$ N. The mV units cancel, leaving N.

\subsection*{Q1(c) \hfill [3 marks]}

For zero-order hold, F = e\textasciicircum{}(AT) and G = integral\_0\textasciicircum{}T e\textasciicircum{}(Aτ)B dτ.
Here:\[\begin{aligned}A = -2, B = 1, T = 0.1 s\end{aligned}\]
Thus:\[\begin{aligned}F = e^(-0.2) = 0.818731\\ G = (1-e^(-0.2))/2 = 0.0906346\end{aligned}\]
Therefore:\[\begin{aligned}x_{k+1} = 0.818731 x_k + 0.0906346 u_k\end{aligned}\]\par\smallskip With $AT=(-2)(0.1)=-0.2$, $F=e^{-0.2}=0.8187308$. For $G$, $(1-F)/2=(1-0.8187308)/2=0.1812692/2=0.0906346$. Thus the input contribution is positive.

\section*{Question 2 \hfill [10 marks]}
\subsection*{Q2(a) \hfill [2 marks]}

Choose x₁ = θ and x₂ = θ̇.
Then ẋ₁ = x₂ and ẋ₂ = [τ - b x₂ - mgℓ sin(x₁)]/J.
At equilibrium:\[\begin{aligned}x_2 = 0 and τ_0 = mg\ell sin \theta_e\end{aligned}\]
Since |sin θ\_e| less than or equal to 1, an equilibrium exists when |τ₀| less than or equal to mgℓ = 2.4525 N m.\par\smallskip At rest $\dot\theta=0$ and $\ddot\theta=0$, so the damping term vanishes. Rearranging gives $\tau_0=mg\ell\sin\theta_e$. Here $mg\ell=1\times9.81\times0.25=2.4525$ N\,m. Since $-1\leq\sin\theta_e\leq1$, the required torque must lie between $-2.4525$ and $+2.4525$ N\,m.

\subsection*{Q2(b) \hfill [2 marks]}

With x = [θ, θ̇]ᵀ, A = [[0,1], [-(mgℓ/J)cos θ\_e, -b/J]], B = [[0],[1/J]].
At θ\_e = 60°, A = [[0,1], [-61.3125, -2.5]] and B = [[0],[50]].\par\smallskip At the equilibrium, the constant torque balances gravity. For a small perturbation, use $\sin(\theta_e+\delta\theta)\approx\sin\theta_e+\cos\theta_e\,\delta\theta$. The constant terms cancel against $\tau_0$; the remaining linear term gives the local state matrix.

\subsection*{Q2(c) \hfill [4 marks]}

At rest, θ̇ = θ̈ = 0, so the torque balance is Kp(θd-θss) = mgℓ sin θss.
With Kp = 50, θd = π/2 and mgℓ = 2.4525, solve 50(π/2-θss) = 2.4525 sin θss.
This gives θss approximately 1.52181 rad = 87.193°, so the error is 90°-87.193° = 2.807°; it is not zero.\par\smallskip Rearranging the steady-state equation gives $50(\pi/2-\theta_{ss})=2.4525\sin\theta_{ss}$. I used radians in the equation. Solving numerically gives $\theta_{ss}\approx1.522\,$rad $\approx87.19^\circ$, so the gap from $90^\circ$ is about $2.81^\circ$.

\subsection*{Q2(d) \hfill [2 marks]}

With gravity compensation, the gravity term cancels the plant's mgℓ sin θ term.
At rest, Kp(θd-θss)=0, so θss=θd and the steady-state error is zero.
The linearised gravity stiffness term is also cancelled, so the closed-loop poles generally change; the stiffness contribution mgℓ cos θd is removed.\par\smallskip The compensation term cancels the nominal gravity torque in the plant model. At steady state the proportional term must then be zero, which requires $\theta_d-\theta_{ss}=0$. This does not mean the transient response is unchanged: removing the gravity derivative changes the linearized stiffness.

\section*{Question 3 \hfill [10 marks]}
\subsection*{Q3(a) \hfill [2 marks]}

\vspace{2mm}

\subsection*{Q3(b) \hfill [2 marks]}

For H(s)=ωc/(s+ωc), ωc=2π(20)=125.664 rad/s and T=0.001 s.
Tustin gives H(z)=α(1+z⁻¹)/[(1+α)+(α-1)z⁻¹], α=ωcT/2=0.0628319.
Thus:\[\begin{aligned}y[k]=0.881765y[k-1]+0.0591175(x_k+x[k-1])\end{aligned}\]
DC gain = (2x0.0591175)/(1-0.881765) approximately 1.000.\par\smallskip $\omega_c=2\pi(20)=125.664$ rad/s. With $T=0.001$ s, $\alpha=\omega_cT/2=125.664(0.001)/2=0.062832$. The recurrence coefficients are approximately $0.881765$ on the previous output and $0.0591175$ on each input sample. Check: $2(0.0591175)/(1-0.881765)\approx1$, so the DC gain is close to one.

\subsection*{Q3(c) \hfill [4 marks]}

With L=[ℓ₁,ℓ₂]ᵀ and C=[1,0], A-LC=[[-ℓ₁,1], [-ℓ₂,-2]].
Its characteristic polynomial is det(sI-(A-LC)) = (s+ℓ₁)(s+2)+ℓ₂ = s²+(ℓ₁+2)s+(2ℓ₁+ℓ₂).
Desired poles at -10 give (s+10)²=s²+20s+100.
Hence ℓ₁+2=20 gives ℓ₁=18, and 2ℓ₁+ℓ₂=100 gives ℓ₂=64.\par\smallskip Expanding $(s+\ell_1)(s+2)+\ell_2$ gives $s^2+(\ell_1+2)s+(2\ell_1+\ell_2)$. The desired polynomial is $(s+10)^2=s^2+20s+100$. Matching the $s$ coefficient gives $\ell_1=18$; matching the constant term then gives $36+\ell_2=100$, hence $\ell_2=64$.

\subsection*{Q3(d) \hfill [2 marks]}

An observer uses the plant model and known control input to estimate unmeasured states such as velocity, and its poles can be designed.
A simple low-pass filter cannot exploit the state-space model in this way.
A filter may be safer when the model is poorly known or unreliable, because it relies less on model accuracy.\par\smallskip The observer uses the plant model and the measured output to reconstruct the unmeasured state. Its convergence depends on the model and the observer poles. A low-pass filter is simpler, but generally introduces phase lag and does not explicitly use the plant dynamics.

\end{document}